\section{Recommendation Process}
\label{sec:recommendation-process}

The recommendation process takes the user's input and produces a 
ranked list of ML methods. The process has two main stages. In 
the first stage, candidate methods are retrieved from the knowledge graph. 
In the second stage, 
if a problem description is provided, the candidates are ranked 
using semantic similarity to relevant research articles. A feedback 
boost is applied as a final adjustment to the ranking.

Figure~\ref{fig:recommendation-sequence} shows the sequence of 
steps in the recommendation process, covering ontology-based 
candidate retrieval and semantic similarity ranking. The steps 
before the conditional block correspond to candidate retrieval, 
described in Section~\ref{sec:ontology-candidate-retrieval}, 
while the steps inside it are the semantic similarity ranking 
steps that run only when a problem description is provided, 
described in Sections~\ref{sec:article-similarity-scoring} 
and~\ref{sec:method-scoring}. Feedback is not included in detail 
in the diagram. It is described separately in 
Section~\ref{sec:feedback-boost}.

\begin{landscape}
\begin{figure}[H]
    \centering
    \includegraphics[width=\linewidth, height=\textheight, keepaspectratio]
        {Figures/06-implementation/recommendation-sequence.png}
    \caption[Recommendation process sequence diagram]{Sequence 
    diagram of the recommendation process in MLguide, showing 
    ontology-based candidate retrieval and semantic similarity 
    ranking.}
    \label{fig:recommendation-sequence}
\end{figure}
\end{landscape}

\subsection{Ontology-Based Candidate Retrieval}
\label{sec:ontology-candidate-retrieval}

The first step is to retrieve candidate methods from the knowledge 
graph using a SPARQL query against GraphDB. The query goes through 
the graph from articles to methods to approaches. It finds articles 
related to the selected ML area, follows the 
\texttt{mentionsMethod} relation to \texttt{mla:Method} instances, 
and then follows \texttt{skos:exactMatch} to the corresponding 
\texttt{ML\_approach} instances. A candidate approach is only 
included if at least one article in the knowledge graph mentions 
a method that maps to it, which grounds the candidate set in the 
research literature.

The properties used as filter options for the user are listed in 
Table~\ref{tab:retrieval-properties}. They were selected because 
they describe characteristics of an ML method that can be matched 
against the user's problem context. ML area is the only hard filter. 
Most articles in the knowledge 
graph are assigned an ML area, as shown in 
Section~\ref{sec:resulting-database}, so restricting to the selected area 
removes few relevant candidates. Articles without an assigned ML 
area are not retrieved when an area is selected, which is 
accepted because the ML area is the most reliable of the available 
properties. The remaining properties have 
incomplete coverage, as shown in 
Table~\ref{tab:ml-approach-coverage-extended}. They are therefore 
not used as hard filters, since excluding methods on the basis of 
incomplete properties would risk removing valid candidates and 
produce a misleading ranking.

% \textbf{Referere til \ref{sec:machine-learning-methods} ? Og den manglende delen om attributter på metoder}

\begin{table}[H]
\centering
\caption[Ontology properties used in candidate retrieval]
{Ontology properties used as filter options in candidate retrieval.}
\label{tab:retrieval-properties}
\begin{tabularx}{\textwidth}{llX}
\toprule
\textbf{Property} & \textbf{Source} & \textbf{Example values} \\
\midrule
\texttt{has\_ml\_area} & \texttt{mla:Article} 
  & \texttt{supervised\_learning}, 
    \texttt{reinforcement\_learning} \\
\midrule
\texttt{used\_for} & \texttt{ML\_approach} 
  & \texttt{classification}, 
    \texttt{regression}, 
    \texttt{anomaly\_detection} \\
\midrule
\texttt{performance} & \texttt{ML\_approach} 
  & \texttt{fast\_training}, 
    \texttt{accurate}, 
    \texttt{explainable} \\
\midrule
\texttt{has\_dataset\_type} & \texttt{ML\_task} 
  via \texttt{used\_for}
  & \texttt{tabular}, 
    \texttt{image}, 
    \texttt{text} \\
\bottomrule
\end{tabularx}
\end{table}

Task, performance, and dataset type are optional filters. 
For each candidate method, the user's selected filter values 
are matched against the method's property values.
Dataset type is accessible indirectly through the tasks linked via 
\texttt{used\_for}, as shown in Table~\ref{tab:retrieval-properties}. 
When no problem description is given, the candidates are sorted by 
the number of matches, with the number of supporting articles as a 
tiebreaker, and the ordered list is returned directly, adjusted 
only for feedback. 
When a problem description is provided, the candidate set is passed on to 
the similarity ranking, which determines the final order. The 
filter matches are then returned as information about each method, 
but do not affect the ranking.

The properties \texttt{possible\_if} and 
\texttt{not\_recommended\_if} are always returned alongside each 
method as informational notes only, and are never used for 
sorting. Their coverage is too low 
(Table~\ref{tab:ml-approach-coverage-extended}) to filter or rank 
on without risking the exclusion of valid candidates.

\subsection{Article Similarity Scoring}
\label{sec:article-similarity-scoring}

When a problem description is provided, articles in the database 
are ranked by their semantic similarity to a problem description query, using the 
two-stage retrieval pipeline combining a bi-encoder and a cross-encoder, as described in 
Section~\ref{sec:cross-encoder-reranking}. Similarity is 
measured using cosine similarity, which is described as the most common metric in
Section~\ref{sec:text-embeddings-and-semantic-similarity}.

In the first stage, the problem description is encoded using 
the bi-encoder
\texttt{all-mpnet-base-v2} \parencite{sbert_mpnet}, a 
sentence-transformer model based on MPNet \parencite{song2020mpnet}. 
The model applies the SBERT approach described in 
Section~\ref{sec:transformer-based-sentence-embedding-models}, 
producing sentence embeddings that can be compared using cosine similarity. This model was used in the pre-master 
project \parencite{bergsto2026mlselection} to compute the article 
embeddings stored in the database, so the same model must be used 
here to ensure the problem description is encoded in the 
same embedding space. The model produces 768-dimensional embeddings and has been shown to 
perform well on abstract-level semantic similarity tasks in 
scientific literature 
\parencite{galli2024performance_pre_trained_sentence_transformer_models, 
colangelo2025comparative_analysis_sentence_transformer_models}. 
Article scores are computed as a weighted sum of cosine 
similarities between the query embedding and the precomputed 
abstract and title embeddings:

\begin{equation}
    \text{score}_a = w_{\text{abs}} \cdot \text{sim}_{\text{abs}}(a)
                   + w_{\text{title}} \cdot \text{sim}_{\text{title}}(a)
    \label{eq:article-score}
\end{equation}

where $w_{\text{abs}} = 0.7$ and $w_{\text{title}} = 0.3$. The 
higher weight on abstracts reflects that abstracts contain more 
information about the method and application context than titles. 
These weights are set as defaults,
but can be adjusted. The top $k = 50$ articles are retained
for the reranking step.

In the second stage, the top-$k$ articles are reranked using 
\texttt{ms-marco-MiniLM-L6-v2} \parencite{sbert_crossencoder_msmarco}, a 
cross-encoder fine-tuned on the MS MARCO passage ranking dataset \parencite{bajaj2018ms_marco_machine_reading_comprehension_dataset}.
The model was chosen because it achieves strong 
ranking performance while being lightweight enough for query-time 
inference over a small candidate set. Among 14 publicly 
available cross-encoder models evaluated on MS MARCO, 
it was found to offer the best performance 
\parencite{thakur2021beir_heterogenous_benchmark_zero_shot_evaluation_ir_models}. 
As described in Section~\ref{sec:cross-encoder-reranking}, a cross-encoder 
processes the query and each article jointly, producing more 
accurate relevance scores than the bi-encoder alone.

\subsection{Method Scoring}
\label{sec:method-scoring}

After article reranking, each candidate method is scored from the 
articles that mention it. Each such article among the top $k$ adds 
to the method's score, and articles ranked higher add more. Using 
rank positions rather than the raw scores from each model gives 
weight to both broad semantic similarity from the bi-encoder and 
precise relevance from the cross-encoder, and avoids differences 
in scale between the two. This follows the principle behind 
Reciprocal Rank Fusion described in 
Section~\ref{sec:rank-fusion}:

\begin{equation}
    \text{score}_m = \sum_{\substack{a \,\in\, \text{top-}k \\ 
    m \,\in\, a}}
        \left(
            \frac{w_b}{r_b(a) + 1} + \frac{w_c}{r_c(a) + 1}
        \right)
\label{eq:method-score}
\end{equation}

where $r_b(a)$ and $r_c(a)$ are the zero-based rank positions of 
article $a$ in the bi-encoder ranking and cross-encoder ranking 
respectively. The weights $w_b$ and $w_c$ are normalised with 
$w_b + w_c = 1$, and both default to $0.5$. The choice to combine 
both signals was made by manually checking output for queries of 
different types and levels of detail, since no ground-truth 
ranking data exists. Using only the cross-encoder score is also 
possible, and was the initial approach. Both the signal weights 
and this choice are configurable.

\subsection{Feedback Boost}
\label{sec:feedback-boost}

After methods are ranked by their combined article-based score, 
user feedback can adjust the final ordering. 
When a user rates a method on a 1--5 scale, 
the rating is stored and used in subsequent recommendations. 
As discussed in Section~\ref{sec:rating-based-user-feedback}, 
rating estimates based on few observations are unreliable. 
A smoothed score is therefore computed using Bayesian smoothing 
towards a neutral prior 
\parencite[Chapter~3.3]{murphy2012machine_learning_probabilistic_perspective}:

\begin{equation}
    \hat{r}_m = \frac{s \cdot \mu_0 + n \cdot \bar{r}_m}{s + n}
    \label{eq:smoothed-rating}
\end{equation}

where $\bar{r}_m$ is the observed mean rating, $n$ is the 
number of ratings, $s = 5$ is the prior strength, and 
$\mu_0 = 3$ is the prior mean on a 1--5 scale. When $n$ is 
small, the estimate is pulled towards $\mu_0$. As $n$ grows, 
the observed mean dominates.
The prior strength of 5 means that the observed ratings and the 
prior carry equal weight at around five ratings, below this the 
prior dominates, and above it the observed mean does. 
This is a low threshold, set with the expectation that the 
system will initially have a small user base. The value is a 
default and can be adjusted.
The smoothed rating is then mapped to $[-1, 1]$ relative to the 
neutral rating of 3 to give the feedback score:
\begin{equation}
    \text{feedback\_score}(m) = \frac{\hat{r}_m - 3}{2}
\label{eq:feedback-score}
\end{equation}

The feedback score is combined with the similarity-based 
ranking using a rank-based score, so that the two signals are 
on a comparable scale:

\begin{equation}
    \text{combined\_score}(m) = \frac{1}{\text{rank}(m) + 1}
                              + w_f \cdot \text{feedback\_score}(m)
    \label{eq:combined-score}
\end{equation}

where $\text{rank}(m)$ is the zero-based position after 
similarity ranking, and $w_f = 0.15$  is an adjustable weight value. 
The low weight 
means that feedback can adjust a method's position by a few 
places, but cannot override the similarity-based ranking 
entirely.

Figure~\ref{fig:ranking-flow} illustrates the full ranking 
process, from initial retrieval of candidate methods, through 
similarity scoring and feedback boost. Arrows are labeled with
the default weight values.  

\begin{figure}[H]
    \centering
    \includegraphics[width=0.8\textwidth]{Figures/06-implementation/ranking-flow.pdf}
    \caption[Ranking flow]{Flowchart of the ranking process in 
    MLguide.}
    \label{fig:ranking-flow}
\end{figure}